\section{Interpretation, integration, and further research}
\label{sec:research}

\subsection{What a positive power separation already follows from}
\label{subsec:composition}

It is useful to isolate a short logical consequence of the prior
arbitrary-factor result. This prevents a quantitative refinement from
being mistaken for the first proof of a qualitative statement.

\begin{proposition}[Power separation from a balanced seed]
\label{prop:composition}
Suppose a Boolean function \(f\) satisfies
\(L_{\mathrm{abs}}(f)>\sqrt{\deg f+1}\). Then there is a balanced
Boolean function \(g\), of some degree \(d>1\), such that
\(L(g)>\sqrt d\). Its iterated block compositions \(g_j\) satisfy
\[
 \deg g_j=d^j,\qquad L(g_j)=L(g)^j
 =(\deg g_j)^{\alpha},\qquad
 \alpha=\frac{\log L(g)}{\log d}>\frac12.
\]
\end{proposition}

\begin{proof}
Changing coordinate signs makes every singleton coefficient of \(f\)
nonnegative. If its mean is negative, replace it by \(-f(-x)\);
this changes the sign of the mean while preserving those nonnegative
singleton coefficients. We may therefore suppose that \(\E f\ge0\)
and \(L(f)=L_{\mathrm{abs}}(f)\).
Define
\[
 g(z,x)=z f(zx).
\]
On a Walsh monomial \(\chi_S(x)\), this transformation gives
\(z^{|S|+1}\chi_S(x)\), whose total degree is odd and at most
\(|S|+1\). Thus \(g\) is odd under simultaneous sign reversal,
is balanced, and has degree \(d\le\deg f+1\).
The old singleton coefficients are unchanged and the new \(z\)
coefficient is \(\E f\). Hence
\[
 L(g)=L(f)+\E f>\sqrt{\deg f+1}\ge\sqrt d.
\]
The linear upper bound \(L(g)\le d\) implies \(d>1\).

Balanced block composition multiplies first-level signed Fourier sums
\cite[Fact 2.7]{BlaisTanWan2015}. To see it directly, expand the outer
function in Walsh characters of its independent, balanced inner
outputs. A singleton input coefficient can only come from a singleton
outer term and a singleton coefficient in that inner block; summing
these products multiplies the two singleton sums.
Degree is multiplicative as well: maximal-degree outer and inner
monomials yield a nonzero monomial on their disjoint blocks, and
different outer supports cannot cancel it. Induction proves the two
identities for \(g_j\) and hence the power law.
\end{proof}

The arbitrary-factor theorem of \cite{OpenAI2026} supplies the
hypothesis of Proposition~\ref{prop:composition}. The normal-transfer
sections of this article instead give one entirely specified route
to a positive exponent, with no stage-dependent existence choice for
the averaging size. Both routes are finite, but their useful outputs
are different: the first is a short qualitative deduction; the second
records explicit construction parameters, moment bounds, and dimension
estimates.

\begin{corollary}[A power lower bound for the retention profile]
\label{cor:profile-power}
Use the fixed-count construction with an admissible gain
\(0<\eta<1\), and put
\[
 r=\frac m{m+1},\qquad q=\frac{m+\eta}{m+1},\qquad
 \alpha=\frac{\log q}{\log r}\in(0,1).
\]
Then the profile from Proposition~\ref{prop:profile} satisfies
\[
 \Phi(\rho)\ge q\rho^\alpha\qquad(0<\rho\le1).
\]
In particular, \(\Phi(\rho)/\rho\to\infty\) as \(\rho\downarrow0\).
\end{corollary}

\begin{proof}
At stage \(k\), the construction has
\[
 N_k=((m+1)M)^k,\qquad
 D_k=(mM)^k=r^kN_k,\qquad
 v(F_k)\ge[M(m+\eta)]^k=q^kN_k.
\]
The supremum formula for \(\Phi\) therefore gives
\(\Phi(r^k)\ge q^k=(r^k)^\alpha\).
If \(0<\rho<1\), choose \(k\ge1\) with
\(r^k\le\rho<r^{k-1}\).
Monotonicity yields
\(\Phi(\rho)\ge(r^k)^\alpha\ge(r\rho)^\alpha
=q\rho^\alpha\). At \(\rho=1\) the assertion follows from
\(\Phi(1)=1\). Since \(\alpha<1\), the quotient diverges.
\end{proof}

The averaging count cancels from this profile bound. Its role is to
make each finite stage valid, while the retained fraction and degree
fraction determine the profile exponent.

\subsection{What the certificates establish}

The proof of the codimension threshold is analytic and exact. Its
finite counterexample is checked in two ways: the displayed selector
formula and direct summation over the \(5^5=3125\) block-sum tuples,
weighted by the exact numbers of sign strings in each tuple.
The grouped calculation is a complete calculation on the twenty-bit
cube because every rule and score is constant on each such tuple.
No sampling is used.

The degree certificate has two independent ingredients. The
cancellation identity gives a pointwise upper bound for every cell.
A nonzero Fourier coefficient gives the matching lower bound for the
exceptional cell. The readout's degree is certified by a separate
nonzero coefficient. The full score law and its second, third, and
fourth moments are recomputed in exact fractions.

The optional exhaustive search uses rational event probabilities and
conditional moments. Its sign tests use integer arithmetic or exact
fractions, so no floating-point tolerance determines whether a gain
is positive. The six-bit, two-leaf row was also computed with an
independently written integer cross-multiplication implementation.
This checks the stated finite search domain; it does not convert a
restricted search into a proof about all partitions.

The explicit Gaussian certificate verifies rational inequalities used
to bound the analytic constants. It does not numerically integrate
tens of millions of tiny intervals or construct the final enormous
Boolean function. The Gaussian density and tail bounds are proved
symbolically in the article. Dyadic constants whose denominators would
have billions of bits are stored by their integer exponents.

\subsection{A proposed ProveIt integration}

The package is a self-contained report directory with a main
\texttt{fourier\_codimension.tex}, its section files, the compiled PDF,
figures, Python certificates, data outputs, and a provenance record.
It can be placed alongside the repository's existing mathematical
reports without importing or modifying the earlier source manuscripts.
The source uses ordinary \LaTeX{} packages and an inline bibliography.
The exact certificates use Python's standard library; only figure
regeneration requires Matplotlib.

A formalization can be developed in layers.
\begin{enumerate}
\item Establish general finite Walsh orthogonality, support
complementation under multiplication by top parity, and degree
bookkeeping for indicators. The inspected declarations
\texttt{sum\_neg\_one\_pow\_binaryWeight\_land} and
\texttt{sum\_thueMorseSign\_mul\_walsh} in the existing
Thue--Morse development provide related finite identities
\cite{ProveItWalsh}; additional abstraction is required.
\item Formalize the interlacing-lattice variance lemma and the
two-point total-variation projection. These are finite sums or
elementary expectations and are independent of asymptotic probability.
\item Derive the cell theorem and its equality and stability
consequences. Careful parity and affine-change conventions make this
layer reusable.
\item Encode the twenty-bit rule by its five block sums. Prove the
indicator identity symbolically, then check the rational moment table
with kernel-verifiable finite arithmetic.
\item Add the finite selector formula and tensorization. The optional
enumeration can later be turned into a reflected decision procedure
or an independently checkable certificate format.
\item Formalize the quantitative amplifier separately, reusing a
Berry--Esseen theorem if an appropriate library result is available.
This layer needs measure-theoretic conditional expectation and a
normal-approximation bound, and is substantially larger.
\end{enumerate}
These are proposed tasks, not names of modules already present.
No Lean build or kernel verification of the new theorems is claimed.

\subsection{Further research questions}

The following problems distinguish structural classification,
quantitative improvement, and formal verification. Several have small
finite first cases, while others concern a genuine asymptotic limit.

\begin{question}[Global minimum dimension]
What is the least \(N\) for which there is any observation \(F\) with
\(v(F)>D(F)\)? Is it twenty? The present minimality theorem only
covers the binomial-block reporting family. A complete search must
allow unrelated cell geometries and enforce that the chosen cells
form a partition.
\end{question}

A promising first step is to search for a single low-degree cell with
small residual variance, then determine whether its complement can
be partitioned at the same degree cost. These two tasks should be
kept separate: a good isolated cell need not admit a useful low-cost
completion.

\begin{question}[The exact codimension-four infimum]
Determine
\[
 a_4=\inf_{\substack{N\ge4,\ C\ne\varnothing\\
                    \deg\1_C\le N-4}}\Var(H_N\mid C).
\]
The results here give \(3\le a_4\le1147/289<4\).
Does the full cube constraint force a gap above three, or can actual
cells approach the abstract alternating-moment infimum?
\end{question}

The lower endpoint is not attained by a cell: equality at three
would force a translated three-sign law, which violates the fourth
alternating cancellation. Any sequence of actual cells approaching
three must therefore have unbounded dimension.

\begin{question}[Higher-codimension moment minima]
For \(k\ge5\), determine
\[
 b_k=\inf\left\{\Var(Y):
  Y\in\Z,\ \E|Y|^{k-1}<\infty,\
  \E[(-1)^YY^j]=0\ (0\le j<k)\right\}.
\]
We prove \(b_1=1/4\), \(b_2=1/2\), and \(b_3=b_4=3/4\),
with nonattainment at \(k=4\). Does \(b_k\) remain bounded,
or must it eventually grow? Which moment extremizers can be realized
as conditional Hamming-weight laws of low-degree cells?
\end{question}

The last realizability condition includes nonradial Walsh constraints
that the one-dimensional moment problem discards. Positivity of a
Hamming-weight distribution is therefore only the beginning of a
cube construction.

\begin{question}[Geometry of equality cells]
Classify the cells of degree at most \(N-k\), for \(k=1,2,3\),
whose conditional sum law is a translated \(k\)-sign sum.
Which are coordinate subcubes, which are not, and what additional
assumptions force a geometric classification?
\end{question}

Natural extra assumptions include invariance under selected
coordinate permutations, monotonicity, bounded numbers of cell
values under restriction, or prescribed affine symmetries.
The distributional equality theorem alone does not choose among
those hypotheses.

\begin{question}[Optimal stability for actual cube cells]
Are the sharp local total-variation moduli in the abstract
alternating-moment class also sharp for uniform cube cells?
Can the constants in the averaged joint-law theorem be improved
using compatibility among the different cells of a partition?
\end{question}

An exact stability optimizer in the moment class need not have
probabilities realizable as integer counts within the available
Hamming layers of a cube cell.
Approximating its weights while preserving all the high-degree
Fourier cancellations is a substantive issue.

\begin{question}[The retained-variance surface]
Determine the extremal function \(V(N,D)\), or its asymptotic
profile \(\Phi\) from Proposition~\ref{prop:profile}.
How close is the tensor lower bound
\[
 D+\frac9{34816}
 \min\left\{\left\lfloor D/16\right\rfloor,
             \left\lfloor(N-D)/4\right\rfloor\right\}
\]
to the optimum? Are there finite seeds with much larger advantage
per bit, per degree, or per exceptional-cell probability?
\end{question}

\begin{question}[Sharper finite reporters]
Optimize the exact selector formula under non-singleton tests,
unequal block sizes, nonidentical input laws, or more than one
exceptional tag. Can one prove finite-dimensional dual bounds
that replace exhaustive enumeration?
\end{question}

The arbitrary-event formula separates a squared mean mismatch from
within-test variances. This suggests convex relaxations for fixed
tests and exact dynamic programming for the selector. Any relaxation
should retain an exact rational certificate for its final candidate.

\begin{question}[A direct discrete amplifier]
Can a small finite score law, perhaps the explicit twenty-bit law,
be iterated with a fixed small number of copies while preserving a
strict multiplicative variance advantage? A proof would avoid the
very large Gaussian grid and might produce an explicit Boolean
counterexample that is practical to tabulate.
\end{question}

The score-law perturbation formula gives a finite state space for
testing such a mechanism. An invariant family of distributions or a
Lyapunov inequality could replace uniform closeness to the normal law.

\begin{question}[Useful numerical exponents]
What explicit exponent \(\varepsilon>0\) can be certified in
\[
 L_{\mathrm{abs}}(f)\ge c(\deg f)^{1/2+\varepsilon}
\]
with moderate construction parameters? Our numerical exponent is
deliberately conservative. The central density-ratio estimate,
interval probabilities, and generic Berry--Esseen transfer each
offer distinct places for improvement.
\end{question}

An optimization should track the resulting dimension and
description length as well as the exponent. A better asymptotic
number can still correspond to a much harder finite object to
construct.

\begin{question}[Limits of first-moment readout]
What additional conditions on an observation guarantee that a
retained-variance excess produces a comparably strong first absolute
moment? Can one classify rare-event perturbations that increase
variance while leaving the majority first moment unchanged, as in
\eqref{eq:witness-score-law}?
\end{question}

This question distinguishes optimizing \(v(F)/D(F)\) from optimizing
the final Boolean Fourier quantity. Bounds using fourth moments are
robust, but may lose information that a sharper distributional
comparison could preserve.

\begin{question}[Weighted coordinate sums]
For \(H_w=\sum_iw_iX_i\), what replaces the codimension threshold and
its equality laws? Coordinate sign changes handle \(w_i\in\{-1,1\}\),
but arbitrary magnitudes destroy the common interlacing lattice.
Are there sharp bounds in terms of the smallest weights, subset sums,
or arithmetic structure of \(w\)?
\end{question}

\begin{question}[Query models and observation degree]
When can a low-degree reporter be implemented by a shallow coordinate
or parity decision tree, and what lower bound on query cost follows
from a retained-variance advantage? Relate the present partitions to
the query and spectral-sample statements already present in ProveIt,
without identifying polynomial degree with query depth.
\end{question}

\begin{question}[A compact formal certificate]
Can the twenty-bit witness, the codimension obstruction, and the
scoped minimality search be packaged into a small Lean development
whose trusted base consists only of the kernel and standard finite
arithmetic? A particularly useful deliverable would check both the
degree upper bound and all exact cell moments from one compact
description of the reporting rule.
\end{question}

The immediate priorities are the global minimum-dimension problem,
the codimension-four infimum, and a practical discrete amplifier.
They directly test whether the structural boundary found here is the
beginning of a broader classification or the first member of a much
larger collection of efficient constructions.
